\chapter{Introduction}\label{ch:introduction}

\section{Motivation}\label{motivation}

Mobile robots usually obtain static stability from a support polygon
formed by three or more ground contacts. A dynamically balanced robot
accepts a narrower support region and maintains the body near an
unstable upright equilibrium by continuously commanding its wheels. This
arrangement can reduce the platform width and enable rapid changes in
the direction of travel, but it also makes stabilization a prerequisite
for every other task. Sensor bias, computation delay, actuator
saturation, gearbox compliance, wheel slip, and power interruption can
therefore affect both mobility and safety. The wheeled inverted pendulum
is a standard model for studying this class of system because it exposes
the coupling between translation and body pitch while remaining
tractable for control design \cite{grasser2002joe}, \cite{pathak2005wip}.

A conventional two-wheel balancing robot can move longitudinally and
turn about a vertical axis, but it cannot command lateral translation
without first changing heading. Mecanum wheels introduce passive rollers
whose axes are inclined relative to the wheel plane. With suitable
roller orientations and independent wheel actuation, the roller
constraints redirect wheel forces so that a mobile base can produce
longitudinal, lateral, and yaw motion \cite{pin1994omnidirectional,ilon1975mecanum,gfrerrer2008geometry,indiveri2009swedish,ferriere1996omnimobile,muir1987kinematic}. Most four-wheel
Mecanum platforms place the wheels at the corners of a rectangular
chassis and rely on a broad, statically stable footprint. Zephyr instead
places four Mecanum wheels on one common wheel-axis line and uses active
pitch stabilization. The concept therefore combines the control problem
of a wheeled inverted pendulum with the kinematic allocation problem of
a redundant omnidirectional drive.

The collinear arrangement is not treated in this thesis as a rectangular
Mecanum base with a vanishing wheelbase. The wheel-center positions,
roller signs, and balance coordinate are defined explicitly, and the
mobility of the actual arrangement is determined from the rank of its
geometry matrix. This treatment is necessary because lateral and yaw
authority depend on the measured axial locations and roller handedness.
An incorrect sign convention could produce an apparently plausible model
that commands the wrong combination of wheels.

\begin{longtable}[]{@{}
  >{\raggedright\arraybackslash}p{(\columnwidth - 0\tabcolsep) * \real{1.0000}}@{}}
\toprule\noalign{}
\begin{minipage}[b]{\linewidth}\raggedright
\textless INSERT FIGURE 1.1: Create an annotated orthographic or
isometric view of Zephyr showing all four wheel centers on the common
axis, wheel numbers 1 through 4, left- and right-handed rollers, body
x-y-z axes, positive pitch and yaw directions, center of mass, and the
plane of ground contact\textgreater{}\\
What the figure should show: Create an annotated orthographic or
isometric view of Zephyr showing all four wheel centers on the common
axis, wheel numbers 1 through 4, left- and right-handed rollers, body
x-y-z axes, positive pitch and yaw directions, center of mass, and the
plane of ground contact\\
Likely source: final SolidWorks assembly or a labeled photograph of the
prototype\strut
\end{minipage} \\
\midrule\noalign{}
\endhead
\bottomrule\noalign{}
\endlastfoot
\end{longtable}

\begin{center}\singlespacing
\phantomsection\addcontentsline{lof}{figure}{\protect\numberline{1.1}Zephyr robot concept and principal degrees of freedom.}
\textbf{Figure 1.1. Zephyr robot concept and principal degrees of freedom.}
\end{center}

\section{Engineering Problem}\label{engineering-problem}

The engineering problem is to determine whether a compact robot with
four collinear Mecanum wheels can be designed and implemented so that it
(1) maintains an upright body, (2) produces controllable motion in the
horizontal plane, and (3) remains electrically and computationally safe
when communication, sensing, or power faults occur. Solving this problem
requires more than selecting a balance gain. The mechanical structure
defines the center-of-mass height and pitch inertia; the roller
arrangement defines the attainable planar wrench; the actuators
introduce torque, speed, backlash, and communication limits; the IMU and
estimator define the observable balance state; and the embedded schedule
determines how much delay and jitter enter the feedback loop.

The project repository identifies Zephyr as a collinear Mecanum-drive
self-balancing robot \cite{zephyr_repo}. The current evidence shows an ESP32-based
bring-up platform, an ISM330DHCX IMU interface, CRSF radio reception, an
ST7789 display interface, CubeMars actuator communication over a 1
Mbit\slash{}s CAN bus, feedback and fault monitoring, and a safety-gated
motor-command path. The four-actuator data structures exist, but the
inspected snapshot enabled only a partial actuator configuration and did
not contain a completed balance-control algorithm, verified wheel-torque
allocation, or hardware performance data. A NUCLEO-H723ZG\slash{}STM32H723ZG
implementation exists as a scaffold or migration target rather than the
demonstrated primary platform. The KiCad design is also a developing
design rather than a manufactured and validated controller board.

\begin{longtable}[]{@{}
  >{\raggedright\arraybackslash}p{(\columnwidth - 0\tabcolsep) * \real{1.0000}}@{}}
\toprule\noalign{}
\begin{minipage}[b]{\linewidth}\raggedright
\textless VERIFY FROM REPOSITORY: Reinspect the current main branch and
all additional branches after exporting the repository as a ZIP. Confirm
the commit hash used for the final thesis, the exact contents of
03\_electrical, 04\_firmware, 05\_simulation, and 99\_archive, and
whether any balance-controller, wheel-allocation, CAD, test-data, or
PCB-routing files were added after commit 97e5a0a.\textgreater{}
\end{minipage} \\
\midrule\noalign{}
\endhead
\bottomrule\noalign{}
\endlastfoot
\end{longtable}

\section{Research Questions and Objectives}\label{research-questions-and-objectives}

The thesis is organized around questions that can be answered by design
evidence, analysis, and repeatable testing. The questions avoid a claim
of novelty based solely on using Mecanum wheels; they focus on the
consequences of combining collinear geometry with dynamic stabilization.

\par\medskip\noindent\singlespacing
\phantomsection\addcontentsline{lot}{table}{\protect\numberline{1.1}Research questions, objectives, and verification evidence}
\textbf{Table 1.1. Research questions, objectives, and verification evidence}\par

\begin{longtable}[]{@{}
  >{\raggedright\arraybackslash}p{(\columnwidth - 4\tabcolsep) * \real{0.3000}}
  >{\raggedright\arraybackslash}p{(\columnwidth - 4\tabcolsep) * \real{0.3500}}
  >{\raggedright\arraybackslash}p{(\columnwidth - 4\tabcolsep) * \real{0.3500}}@{}}
\toprule\noalign{}
\begin{minipage}[b]{\linewidth}\raggedright
\textbf{Research question}
\end{minipage} & \begin{minipage}[b]{\linewidth}\raggedright
\textbf{Engineering objective}
\end{minipage} & \begin{minipage}[b]{\linewidth}\raggedright
\textbf{Required verification evidence}
\end{minipage} \\
\midrule\noalign{}
\endhead
\bottomrule\noalign{}
\endlastfoot
RQ1. What planar motions are kinematically achievable with the measured
collinear wheel geometry and roller pattern? & Derive the geometry
matrix, determine its rank, and establish wheel-speed and
wrench-allocation mappings. & Measured wheel coordinates and roller
signs; analytical rank test; numerical inverse\slash{}forward kinematic
tests. \\
RQ2. Can the completed robot maintain upright balance within defined
limits while stationary and while receiving motion commands? & Implement
a discrete balance controller with verified state estimation,
saturation, and fault handling. & Pitch, pitch-rate, command, current,
and timing logs from repeated balance and motion tests. \\
RQ3. How do actuator limits, roller slip, delay, and mass-property
uncertainty affect performance? & Develop symbolic and numerical models
and perform sensitivity analysis before hardware tuning. & Parameter
record, simulation cases, model-to-hardware comparison, and uncertainty
bounds. \\
RQ4. Does the embedded and electrical architecture support safe,
maintainable experimentation? & Separate portable logic from hardware
code, enforce command ownership, and verify power and CAN integrity. &
Code-path review, task-timing measurements, power checks, CAN error
statistics, and fault-injection results. \\
\end{longtable}

\emph{Note.} Quantitative acceptance thresholds remain proposed until
the author and advisor approve final performance requirements.

The primary objective is to produce an integrated research platform and
a defensible method for evaluating it. Supporting objectives are to
document the mechanical and electrical interfaces, derive a model that
respects the nonstandard wheel arrangement, implement a modular
real-time firmware architecture, establish a safe integration procedure,
and quantify balance and planar-motion performance. The thesis does not
assume that all objectives have already been achieved. Each result is
assigned a status and is linked to evidence or an explicit placeholder.

\section{Scope, Constraints, and Contributions}\label{scope-constraints-and-contributions}

The scope includes mechanical arrangement, embedded electronics,
actuator communication, state estimation, kinematics, reduced dynamics,
control design, simulation, integration, and laboratory validation.
Autonomous navigation, perception, mapping, and application-specific
payloads are outside the present scope. The model treats contact with a
level rigid surface and initially represents each Mecanum wheel by an
idealized roller constraint. Detailed roller compliance and
contact-force variation are reserved for model refinement because
published work shows that individual roller contact and slip can create
motion errors that an ideal kinematic model does not capture
\cite{bayar2020contact}--\cite{barcenas2025robust}.

Several constraints materially shape the design. The four actuators are
integrated CubeMars AK40-10 units with internal field-oriented control,
a 10:1 gearbox, a single internal encoder, and CAN communication
\cite{cubemars2026ak4010}. Their backlash and internal control mode affect low-speed
reversals. The power source is a six-cell lithium-polymer battery, for
which the fully charged voltage must be used when rating converters and
protection components. The current prototype uses an ESP32, while the
STM32H723ZG remains a planned higher-performance target. The final wheel
order, roller pattern, center of mass, inertias, controller gains,
current limits, and sampling rates have not yet been documented with
sufficient evidence for a final thesis and are therefore represented by
explicit verification items.

The supported contribution at this stage is an integrated design
framework rather than a completed performance claim. The framework
combines: a geometry-parameterized Jacobian for four collinear Mecanum
wheels; a reduced nonlinear and linear balance model; a proposed
allocation and control structure; a modular firmware architecture with
safety-owned actuator permission; an electrical design review tied to
the active firmware pin contract; and a test matrix that converts
remaining work into measurable thesis evidence. The originality of the
physical arrangement may be discussed only after a broader prior-art
search and advisor review.

\begin{longtable}[]{@{}
  >{\raggedright\arraybackslash}p{(\columnwidth - 0\tabcolsep) * \real{1.0000}}@{}}
\toprule\noalign{}
\begin{minipage}[b]{\linewidth}\raggedright
\textless INSERT FIGURE 1.2: Show the iterative workflow linking
requirements, CAD and mass properties, schematic and PCB design,
symbolic modeling, controller simulation, firmware bring-up,
current-limited bench tests, restrained balancing, free balancing, and
model-to-hardware comparison\textgreater{}\\
What the figure should show: Show the iterative workflow linking
requirements, CAD and mass properties, schematic and PCB design,
symbolic modeling, controller simulation, firmware bring-up,
current-limited bench tests, restrained balancing, free balancing, and
model-to-hardware comparison\\
Likely source: author-created vector diagram based on the documented
thesis process\strut
\end{minipage} \\
\midrule\noalign{}
\endhead
\bottomrule\noalign{}
\endlastfoot
\end{longtable}

\begin{center}\singlespacing
\phantomsection\addcontentsline{lof}{figure}{\protect\numberline{1.2}Research and development workflow used for the Zephyr project.}
\textbf{Figure 1.2. Research and development workflow used for the Zephyr project.}
\end{center}

\section{Thesis Organization}\label{thesis-organization}

Chapter 2 reviews wheeled inverted pendulums, omnidirectional mobility,
Mecanum-wheel modeling, state estimation, control, and embedded
implementation. Chapter 3 converts the research questions into system
requirements and interfaces. Chapters 4 and 5 document the mechanical
and electrical designs. Chapter 6 derives the kinematic and dynamic
models, and Chapter 7 defines the control and estimation strategy.
Chapter 8 describes the firmware architecture and its demonstrated
status. Chapter 9 specifies the simulation and analysis workflow.
Chapter 10 presents completed and proposed integration procedures, while
Chapter 11 provides ready-to-fill results tables and a cautious
discussion of the evidence available at this draft stage. Chapter 12
summarizes supported conclusions and remaining work. The appendices
preserve parameter, wiring, configuration, test, evidence, and
compliance records needed to complete the final thesis.
